%------------------------------------------------------------------------------%

\usepackage{calculo_varias_variaveis-1}

\title {Multiplicadores de Lagrange}

%------------------------------------------------------------------------------%
\begin{document}

\addtitle

%------------------------------------------------------------------------------%
\section{Extremos Condicionados}

%------------------------------------------------------------------------------%
\begin{frame}{Extremos Condicionados}
  Encontre o máximo (mínimo) da função \(f(x, y, z)\)\\[2mm]
  sujeito a \(g(x, y, z) = 0\)

  \pause
  \vspace{2\baselineskip}
  \emph{$f$} função objetivo

  \pause
  \emph{$g$} restrição
\end{frame}

%------------------------------------------------------------------------------%
\framedgraphic{Extremos Condicionados}{campo_gradiente}{}

%------------------------------------------------------------------------------%
\section{Extremos Locais Restritos a uma Curva}

%------------------------------------------------------------------------------%
\begin{frame}{Teorema do Gradiente Ortogonal}
  Seja \(f(x, y, z)\) uma função diferenciável em uma região

  \pause
  cujo interior contém a curva
  \[
    C\colon r(t)
    = g(t)\veci
    + h(t)\vecj
    + k(t)\veck
  \]

  \pause
  Se $P_0$ é um ponto em $C$ onde $f$ possua um máximo ou mínimo local relativo
  aos seus valores em $C$,
  \pause
  então \(\nabla f\) é ortogonal a $C$ em $P_0$
\end{frame}

%------------------------------------------------------------------------------%
\begin{frame}{Demonstração}
  Os valores de $f$ em $C$ são dados pela composta
  \[
    p(t) = f(g(t), h(t), k(t))
  \]
  \pause
  então
  \[
    \frac{dp}{dt}
    \pause
    = \D{f}{g}\frac{dg}{dt}
    + \D{f}{h}\frac{dh}{dt}
    + \D{f}{k}\frac{dk}{dt}
    \pause
    = \nabla f \cdot \frac{dr}{dt}
  \]

  \pause
  Se $f$ possui máximo (mínimo) em \((a, b, c)\) correspondente a $t_0$
  \[
    \nabla f\evalat{(a, b, c)}{} \cdot \frac{dr}{dt}\evalat{t_0}{} = 0
  \]
\end{frame}

%------------------------------------------------------------------------------%
\section{Multiplicadores de Lagrange}

\framedgraphic{Extremos Condicionados}{LagrangeMultipliers2D}{}

%------------------------------------------------------------------------------%
\begin{frame}{Método dos Multiplicadores de Lagrange}
  Suponha que \(f(x, y, z)\) e \(g(x, y, z)\) sejam diferenciáveis\\[2mm]
  e \(\nabla g\neq 0\) quando \(g(x, y, z)=0\)

  \pause
  \vspace{\baselineskip}
  Para encontrar os valores extremos locais de \(f(x, y, z)\) \\[2mm]
  \pause
  sujeito a restrição \(g(x, y, z)=0\) \\[2mm]
  \pause
  encontre $x$, $y$, $z$ e $\lambda$ tais que
  \[
    \nabla f = \lambda \nabla g
    \qquad\qquad
    g(x, y, z) = 0
  \]
  \pause
  \(\lambda\) é o Multiplicador de Lagrande
\end{frame}

%------------------------------------------------------------------------------%
\addtocounter{exemplo}{1}
\begin{frame}{Exemplo \theexemplo}
  Encontre o maior e menor valores que a função
  \[
    f(x, y) = xy
  \]
  assume na elipse
  \[
    \frac{x^2}{8} + \frac{y^2}{2} = 1
  \]
\end{frame}

\begin{frame}{Exemplo \theexemplo}
  Queremos encontrar os valores extremos da função
  \[
    f(x, y) = xy
  \]
  \pause
  sujeito a restrição
  \[
    g(x, y) = \frac{x^2}{8} + \frac{y^2}{2} - 1 = 0
  \]

  \pause
  Para isso precisamos encontrar os valores $x$, $y$, $z$ e $\lambda$ tais que
  \[
    \nabla f = \lambda \nabla g
    \qquad\qquad
    g(x, y) = 0
  \]
\end{frame}

\begin{frame}{Exemplo \theexemplo}
\begin{columns}
\column{0.4\linewidth}
  Gradiente de $f$
  \[
    \pause
    \D{f}{x}
    \pause
    = \D{}{x}\left(xy\right)
    \pause
    = y
  \]
  \[
    \pause
    \D{f}{y}
    \pause
    = \D{}{y}\left(xy\right)
    \pause
    = x
  \]
  \pause
  \[
    {\color<19>{blue}\nabla f = \vetor{y}{x}}
  \]
\column{0.6\linewidth}
  \pause
  Gradiente de $g$
  \[
    \pause
    \D{g}{x}
    \pause
    = \D{}{x}\left(\frac{x^2}{8} + \frac{y^2}{2} - 1\right)
    \pause
    = \frac{2x}{8}
    \pause
    = \frac{x}{4}
  \]
  \[
    \pause
    \D{g}{y}
    \pause
    = \D{}{y}\left(\frac{x^2}{8} + \frac{y^2}{2} - 1 \right)
    \pause
    = \frac{2y}{2}
    \pause
    = y
  \]
  \pause
  \[
    {\color<19>{blue}\nabla g = \vetor{\dfrac{x}{4}}{y}}
  \]
\end{columns}
\end{frame}

\begin{frame}[t]{Exemplo \theexemplo}
\begin{columns}
\column{0.5\linewidth}
  A equação
  \[
    \nabla f = \lambda \nabla g
  \]
  \pause
  se torna
  \[
    \vetor{y}{x} = \lambda \vetor{\sfrac{x}{4}}{y}
  \]
  {\color<8>{blue}%
  \pause
  \[
    \begin{cases}
      y = \dfrac{\lambda x}{4} \\[2mm]
      x = \lambda y
    \end{cases}
  \]}
\column{0.5\linewidth}
  \pause
  A equação
  \[
    g(x, y) = 0
  \]
  \pause
  se torna
  \[
    \frac{x^2}{8} + \frac{y^2}{2} - 1 = 0
  \]
  \[
    \pause
    \frac{x^2}{8} + \frac{y^2}{2} = 1
  \]
  {\color<8>{blue}%
  \pause
  \[
    x^2 + 4y^2 = 8
  \]}
\end{columns}
\end{frame}

\begin{frame}{Exemplo \theexemplo}
\begin{columns}
\column{0.5\linewidth}
  Precisamos resolver o sistema
  \[
    \begin{cases}
      y = \dfrac{\lambda x}{4} \\
      x = \lambda y \\
      x^2 + 4y^2 = 8
    \end{cases}
  \]
  \pause
  Por substituição
  \[
    y = \dfrac{\lambda}{4} x
    \pause
    = \dfrac{\lambda}{4} \lambda y
    \pause
    = \dfrac{\lambda^2}{4} y
  \]
\column{0.5\linewidth}
  \pause
  Assim {\color<10>{blue}$y=0$}
  \pause
  ou
  \[
    1 = \dfrac{\lambda^2}{4}
  \]
  \[
    \pause
    \lambda^2 = 4
  \]
  \[
    \pause
    \abs{\lambda} = 2
  \]
  \[
    \pause
    {\color<10>{blue}\lambda = \pm 2}
  \]
\end{columns}
\end{frame}

\begin{frame}{Exemplo \theexemplo}
  Caso 1
  \[
    y = 0
  \]

  \pause
  \[
    x = \lambda y
    \pause
    = 0
  \]

  \pause
  \[
    x^2 + 4y^2
    \pause
    = 0
    \pause
    \neq 8
  \]

  \pause
  Não resolve o sistema

  \pause
  \vspace{\baselineskip}
  Portanto \(y\neq 0\)
\end{frame}

\begin{frame}{Exemplo \theexemplo}
  Caso 2

  \pause
  \[
    y \neq 0
    \qquad
    \lambda = \pm 2
  \]

  \pause
  \[
    x
    = \lambda y
    \pause
    = \pm 2 y
  \]
  \[
    \pause
    x^2 = 4y^2
  \]
  \[
    \pause
    x^2 + 4y^2
    \pause
    = 4y^2 + 4y^2
    \pause
    = 8y^2
    \pause
    = 8
  \]

  \pause
  Portanto \(y=\pm 1\)

  \pause
  \vspace{\baselineskip}
  Temos então os pontos
  \((-2, -1)\),
  \((2, -1)\),
  \((-2, 1)\) e
  \((2, 1)\)

\end{frame}

\begin{frame}{Exemplo \theexemplo}
  \onslide<+->{Avaliando a função \(f(x, y)=xy\) nos pontos}
  \begin{align*}
    \onslide<+->{f(           -2,            -1) & = (-2)(-1) = 2 \\[2mm]}
    \onslide<+->{f(\hphantom{-}2,            -1) & = 2(-1) = -2 \\[2mm]}
    \onslide<+->{f(           -2, \hphantom{-}1) & = -2 \times 1 = -2 \\[2mm]}
    \onslide<+->{f(\hphantom{-}2, \hphantom{-}1) & = 2\times 1 = 2}
  \end{align*}

  \onslide<+->{O valor máximo de $f$ é \hspace{1ex} 2 e ocorre nos pontos \((-2, -1)\) e \(( 2, 1)\)}

  \onslide<+->{O valor mínimo de $f$ é $-2$ e ocorre nos pontos \((-2, 1)\) e \((2, -1)\)}
\end{frame}

%------------------------------------------------------------------------------%
\addtocounter{exemplo}{1}
\begin{frame}{Exemplo \theexemplo}
  Encontre os valores máximo e mínimo da função
  \[
    f(x, y) = 3x + 4y
  \]
  na circunferência
  \[
    x^2 + y^2 = 1
  \]
\end{frame}

\begin{frame}{Exemplo \theexemplo}
  Queremos encontrar os valores extremos da função
  \[
    f(x, y) = 3x + 4y
  \]
  \pause
  sujeito a restrição
  \[
    g(x, y) = x^2 + y^2 - 1 = 0
  \]

  \pause
  Para isso precisamos encontrar os valores $x$, $y$, $z$ e $\lambda$ tais que
  \[
    \nabla f = \lambda \nabla g
    \qquad\qquad
    g(x, y) = 0
  \]
\end{frame}

\begin{frame}{Exemplo \theexemplo}
\begin{columns}
\column{0.4\linewidth}
  Gradiente de $f$
  \[
    \pause
    \D{f}{x}
    \pause
    = \D{}{x}\left(3x + 4y\right)
    \pause
    = 3
  \]
  \[
    \pause
    \D{f}{y}
    \pause
    = \D{}{y}\left(3x + 4y\right)
    \pause
    = 4
  \]
  \pause
  \[
    {\color<17>{blue}\nabla f = \vetor{3}{4}}
  \]
\column{0.6\linewidth}
  \pause
  Gradiente de $g$
  \[
    \pause
    \D{g}{x}
    \pause
    = \D{}{x}\left(x^2 + y^2 - 1\right)
    \pause
    = 2x
  \]
  \[
    \pause
    \D{g}{y}
    \pause
    = \D{}{y}\left(x^2 + y^2 - 1\right)
    \pause
    = 2y
  \]
  \pause
  \[
    {\color<17>{blue}\nabla g = \vetor{2x}{2y}}
  \]
\end{columns}
\end{frame}

\begin{frame}[t]{Exemplo \theexemplo}
\begin{columns}
\column{0.5\linewidth}
  A equação
  \[
    \nabla f = \lambda \nabla g
  \]
  \pause
  se torna
  \[
    \vetor{3}{4} = \lambda \vetor{2x}{2y}
  \]
  {\color<7>{blue}%
    \pause
    \[
    \begin{cases}
      x = \dfrac{3}{2\lambda} \\[4mm]
      y = \dfrac{2}{\lambda}
    \end{cases}
  \]}
\column{0.5\linewidth}
  \pause
  A equação
  \[
    g(x, y) = 0
  \]
  \pause
  se torna
  \[
    x^2 + y^2 - 1 = 0
  \]
  {\color<7>{blue}%
    \pause
  \[
    x^2 + y^2 = 1
  \]}
\end{columns}
\end{frame}

\begin{frame}{Exemplo \theexemplo}
\begin{columns}
\column{0.5\linewidth}
  \onslide<+->{Precisamos resolver o sistema
  \[
  \begin{cases}
    x = \dfrac{3}{2\lambda} \\[3mm]
    y = \dfrac{2}{\lambda} \\[3mm]
    x^2 + y^2 = 1
  \end{cases}
  \]}
\column{0.5\linewidth}
  \begin{align*}
    \onslide<+->{x^2 + y^2 & = 1 \\[2mm]}
    \onslide<+->{\left(\frac{3}{2\lambda}\right)^2 +
    \left(\frac{2}{\lambda}\right)^2 & = 1  \\[2mm]}
    \onslide<+->{\frac{9}{4\lambda^2} +
    \frac{4}{\lambda^2} & = 1  \\[2mm]}
    \onslide<+->{\frac{9+16}{4\lambda^2} & = 1  \\[2mm]
    4\lambda^2 & = 25  \\[2mm]}
    \onslide<+->{\lambda & = \pm \frac{5}{2}}
  \end{align*}
\end{columns}
\end{frame}

\begin{frame}{Exemplo \theexemplo}
\begin{columns}
\column{0.5\linewidth}
  Quando \(\lambda=\sfrac{5}{2}\)
  \[
    \pause
    x
    = \frac{3}{2\lambda}
    \pause
    = \frac{3}{2\sfrac{5}{2}}
    \pause
    = \frac{3}{5}
  \]
  \[
    y
    = \frac{2}{\lambda}
    \pause
    = \frac{2}{\sfrac{5}{2}}
    \pause
    = \frac{4}{5}
  \]
\column{0.5\linewidth}
  \pause
  Quando \(\lambda=-\sfrac{5}{2}\)
  \[
    \pause
    x
    = \frac{3}{2\lambda}
    \pause
    = -\frac{3}{2\sfrac{5}{2}}
    \pause
    = -\frac{3}{5}
  \]
  \[
    \pause
    y
    = \frac{2}{\lambda}
    \pause
    = -\frac{2}{\sfrac{5}{2}}
    \pause
    = -\frac{4}{5}
  \]
\end{columns}
\end{frame}

\begin{frame}{Exemplo \theexemplo}
  \onslide<+->{Avaliando a função \(f(x, y)=3x + 4y\) nos pontos}
  \begin{align*}
    \onslide<+->{f\left(           -\frac{3}{5},            -\frac{4}{5}\right) &= }
    \onslide<+->{3\left(-\frac{3}{5}\right) + 4\left(-\frac{4}{5}\right) = }
    \onslide<+->{\frac{-9-16}{5} = }
    \onslide<+->{-5 \\[5mm]}
    \onslide<+->{f\left(\hphantom{-}\frac{3}{5}, \hphantom{-}\frac{4}{5}\right) &= }
    \onslide<+->{3\frac{3}{5} + 4\frac{4}{5} = }
    \onslide<+->{\frac{9+16}{5} = }
    \onslide<+->{5}
  \end{align*}

  \onslide<+->{O valor máximo de $f$ é \hspace{1ex} 5 e ocorre no ponto
  \(\left(\frac{3}{5}, \frac{4}{5}\right)\)}

  \vspace{0.3\baselineskip}
  \onslide<+->{O valor mínimo de $f$ é $-5$ e ocorre no ponto
  \(\left(-\frac{3}{5}, -\frac{4}{5}\right)\)}
\end{frame}

%------------------------------------------------------------------------------%
\section{Lista Mínima}

%------------------------------------------------------------------------------%
\begin{frame}{Lista Mínima}
  Cálculo Vol. 2 do Thomas 12\textsuperscript{a} ed. -- Seção 14.8
  \begin{enumerate}
    \item Estudar o texto da seção
    \item Resolver os exercícios: 3, 5, 7, 12, 16, 18, 24 e 30
  \end{enumerate}

  \vfill
  Atenção: A prova é baseada no livro, não nas apresentações
\end{frame}

%------------------------------------------------------------------------------%
\end{document}
%------------------------------------------------------------------------------%
